% !TeX root = ../main.tex
\chapter{Introducción}

Texto de la introducción. Ejemplos de uso de la plantilla:

\begin{itemize}[label=$\bullet$]
	\item Cita entre paréntesis: \citapendiente{}.
	\item Cita en el texto: según \citapendiente{}, \dots
	\item Término del glosario: una \gls{api} expone varios \glspl{endpoint}.
	\item Acrónimo (la primera vez sale completo): \gls{tfg}, y después \gls{tfg}.
	\item Nota a pie con enlace\footnote{\url{https://www.ejemplo.com}}.
	\item Nota pendiente en el margen.\todo{Revisar este párrafo}
\end{itemize}

% ---------- Figura ----------
\begin{figure}[H]
	\centering
	\includegraphics[width=0.6\linewidth]{img/ejemplo.png}
	\caption[Figura de ejemplo]{Figura de ejemplo. (Fuente propia)}
	\label{fig:ejemplo}
\end{figure}

Como se ve en la Figura~\ref{fig:ejemplo}, \dots

% ---------- Tabla ----------
\begin{table}[H]
	\centering
	\renewcommand{\arraystretch}{1.5}
	\begin{tabularx}{\textwidth}{p{3cm} >{\raggedright\arraybackslash}X}
		\textbf{Columna 1} & \textbf{Columna 2} \\
		\midrule
		Dato & Descripción más larga que se ajusta al ancho disponible. \\
		\midrule
		Dato & Otra descripción. \\
	\end{tabularx}
	\caption{Tabla de ejemplo}
	\label{tab:ejemplo}
\end{table}

% ---------- Código ----------
\begin{lstlisting}[language=JavaScript, style=codigo, caption={Código de ejemplo}, label={lst:ejemplo}]
// Comentario de ejemplo
export async function saludar(nombre) {
	const mensaje = "Hola, " + nombre;
	return mensaje;
}
\end{lstlisting}
